\section{Proof of Theorem~\ref{thm:main}}
\label{supp:proof}

Let $\Pi$ satisfy S1 to S4, let $F$ be the pseudorandom function used to derive record tags and
the per-flush random choices, and let $\mathsf{AE}$ be the authenticated encryption scheme. Let
$\mathcal{A}$ be a probabilistic polynomial-time adversary in $\mathbf{G}^{\mathrm{H}}$ that
submits two equal-length histories --- admissible or not --- and issues polynomially many post-$\tau$
operations. The proof proceeds through three games and then exhibits the simulator of
\cref{def:dualepoch}.

\paragraph{Cryptographic model} $\mathsf{AE}$ is
XChaCha20-Poly1305, used as an authenticated cipher that is IND\$-CPA (ciphertexts are
indistinguishable from uniform strings of the same length) and INT-CTXT. Nonces are $192$-bit and
drawn uniformly at random per encryption; the width makes a collision negligible ($q^2/2^{193}$ over
$q$ encryptions) and removes any nonce counter from disk. A header's associated data is a
label naming the engine and format version together with the header's position (region and slot in
Mnemos-P, absolute slot in Mnemos-M), so a header moved to another slot fails authentication; a
payload's associated data is its record's identifier tag, which binds it to its header. Headers,
payloads, the manifest, the table and the journal are sealed under distinct derived keys. The manifest and the partition
table are themselves $\mathsf{AE}$ ciphertexts of fixed length under their own derived keys; the
partition table is a ciphertext, not a PRF image, and is included in the snapshot as one of the
fixed-length uniform strings of game $\mathrm{G}_2$. All working keys are derived from the master by
HKDF, treated as a PRF with independent labels per key, so the derived keys are independent. The
simulator is given the public structural parameters $(B, C, d, \text{block size}, w, \Delta, k)$ and
the leakage $\mathcal{L}$; it is \emph{not} given $b$, the histories, or any plaintext. Adaptivity is
permitted on physical locations the adversary has already observed, which the next paragraph treats.

\paragraph{On adaptivity} $\mathcal{A}$ chooses each post-$\tau$ operation after seeing the leakage
of every earlier one. This does not disturb the argument. Each hybrid step below replaces a
primitive with an idealized version, uniformly over the whole game, rather than fixing the
operation sequence in advance; the reduction answers $\mathcal{A}$'s queries as they arrive and its
advantage is unchanged. The final step shows the view is a function of the leakage and public
randomness alone, which holds pointwise for every sequence $\mathcal{A}$ might produce and
therefore holds for an adaptively chosen one.

\paragraph{Game $\mathrm{G}_0$} The real game $\mathbf{G}^{\mathrm{H}}_{\Pi,\mathcal{A}}$.

\paragraph{Game $\mathrm{G}_1$} As $\mathrm{G}_0$, except that every evaluation of $F$ under a
store key is replaced by a lazily sampled random function. This affects the record tags, the
keyed logical-to-physical partition permutation, the free-slot choice within a rewritten region,
the choice of uniformly random padding regions in a flush, and the per-call probe randomization.
A distinguisher between $\mathrm{G}_0$ and $\mathrm{G}_1$ yields a distinguisher against $F$, so
$|\Pr[\mathrm{G}_0 \Rightarrow 1] - \Pr[\mathrm{G}_1 \Rightarrow 1]| \le
\mathsf{Adv}^{\mathrm{prf}}_{F}(\lambda)$ up to a factor for the number of derived keys.

\paragraph{Game $\mathrm{G}_2$} As $\mathrm{G}_1$, except that every ciphertext written to disk is
replaced by a uniformly random string of the same length. This includes each slot's header and
payload ciphertexts, the manifest, and the partition table, and every re-encryption performed when
a region is rewritten during a flush. The challenger keeps the plaintexts of live records in its
own state, so no stored ciphertext is decrypted. Each ciphertext is produced under a fresh
uniformly random nonce and the number of encryptions is polynomial, so a reduction answering every
encryption with an IND\$-CPA challenge oracle gives
$|\Pr[\mathrm{G}_1 \Rightarrow 1] - \Pr[\mathrm{G}_2 \Rightarrow 1]| \le
\mathsf{Adv}^{\mathrm{ind\$}}_{\mathsf{AE}}(\lambda)$.

\paragraph{Independence of the view in $\mathrm{G}_2$ from $b$} The view consists of the snapshot
at $\tau$ and the leakage of each post-$\tau$ operation. Each component is considered in turn.

\begin{enumerate}
\item \emph{File lengths.} By S1 every file has a length determined by the public parameters, so
the snapshot's shape is independent of $b$.
\item \emph{Snapshot bytes.} By S2 and the change made in $\mathrm{G}_2$, every slot of the
partition file, whether it holds a live record or is free, is a uniformly random string of the
slot length, and the manifest and table are uniformly random strings of fixed length. The snapshot
is therefore a tuple of independent uniform strings whose count and lengths are public parameters,
identically distributed under $b=0$ and $b=1$.
\item \emph{Flush leakage.} By S3 a flush writes exactly $w+1$ whole regions. The regions written
are: one per staged record, being that record's physical region; uniformly random regions to reach
$w$; and one region fixed by a counter that depends only on the flush index. A staged record's
physical region is the relabeling of the nearest-centroid partition of its embedding, a function of
the embedding, the partition table and the relabeling, all independent of $b$ by S4 and
$\mathrm{G}_1$. The uniform and counter choices are independent of $b$. The written bytes are
uniform by item 2. So a flush reveals only a multiset of region indices with a distribution
independent of $b$.
\item \emph{Deletion and update leakage.} By S3 a deletion writes nothing at the time of the call
and records a tombstone in the manifest, which is rewritten wholesale every flush regardless of
content; an update is a deletion followed by a fresh write. The physical removal of a tombstoned
record occurs only when the counter-selected region of a later flush is rewritten, which is already
covered by item 3 and whose region choice is content-independent. Hence deleting or updating a
cohort-H record adds nothing to the view that depends on which region held it.
\item \emph{Search leakage.} A search reads the regions its probe set names. The probe set is a
function of the query, the partition table, the relabeling and fresh per-call randomness, none of
which depends on $b$. Ranking and payload decryption occur in memory over the region images already
read, so no further disk access is issued.
\item \emph{Refusals.} No admissibility assumption is needed: neither history is required to stay
below capacity, and the case is closed by S3's padded refusal. Consider a post-$\tau$
$\mathsf{put}$ whose target partition is full under one history and not under the other. Under both,
S3 requires the store to rewrite the intended region: an accepted write re-seals the region with the
new record among its slots, and a refused write re-seals the same region unchanged. By item~2 every
slot of a rewritten region is a uniform string of the slot length whether or not it holds a live
record, so the two rewrites are identically distributed, and the region index written is the same
function of the (relabeled) target partition in both cases. The observation extends to a
$\mathsf{put}$ that overflows a partition \emph{within} $\sigma_b$, before $\tau$: the refused writes
rewrite their regions, and by item~2 the snapshot does not reveal a region's occupancy, so a history
that overflows a partition and one that does not produce identically distributed snapshots. Hence
the refusal rule is exercised freely under both histories without distinguishing them, and no
admissibility restriction is needed. What overflow costs is availability, not confidentiality: a
padded refusal discards the caller's write. \Cref{prop:capacity} states this as the
capacity--confidentiality trade-off --- padded refusal (chosen here) hides the load at the price of
lost writes, silent refusal keeps the writes at the price of the load --- and \cref{sec:eval-refusal}
measures both, including the utilization at which the losses begin.
\end{enumerate}

Every component is thus a function of the public parameters, the leakage of the operations issued,
and randomness independent of $b$. A simulator given only $\mathcal{L}$ reproduces the view: it
samples the relabeling $\pi$ itself, emits fixed-length uniform strings for the snapshot, and for
each operation emits the region indices $\mathcal{L}^{\mathrm{op}}$ names under $\pi$, drawing the
uniform padding and per-call probe randomness itself. Hence the real and simulated views are
indistinguishable and $\Pr[b'=b] = \tfrac{1}{2}$ in $\mathrm{G}_2$, giving
\[
\AdvH_{\Pi,\mathcal{A}}(\lambda) \le \mathsf{Adv}^{\mathrm{prf}}_{F}(\lambda) +
\mathsf{Adv}^{\mathrm{ind\$}}_{\mathsf{AE}}(\lambda) \le \negl(\lambda).
\]
\hfill$\blacksquare$

\paragraph{Predicate hiding (\cref{cor:predicate})} Let $\mathcal{D}$ be a distribution over
cohort-H histories, let $\phi$ be a predicate on cohort-H contents, and let $\mathcal{P}$ be the
best predictor of $\phi$ that does not see the adversary's view. Suppose an adversary $\mathcal{B}$
predicts $\phi$ from its view with advantage $\delta$ over $\mathcal{P}$, where the advantage is
taken in expectation over $\mathcal{D}$.

Construct a distinguisher $\mathcal{A}$ for $\mathbf{G}^{\mathrm{H}}$ as follows. $\mathcal{A}$
samples $\sigma_0 \leftarrow \mathcal{D} \mid \phi = 0$ and $\sigma_1 \leftarrow \mathcal{D} \mid
\phi = 1$, submits them, runs $\mathcal{B}$ on the resulting view, and outputs $\mathcal{B}$'s
prediction of $\phi$ as its guess for $b$. Conditioned on the sampled pair, $\mathcal{A}$ is correct
exactly when $\mathcal{B}$ predicts $\phi$ correctly, so
$\Pr[b' = b] = \tfrac{1}{2} + \tfrac{1}{2}\,\delta_{\sigma_0,\sigma_1}$ where
$\delta_{\sigma_0,\sigma_1}$ is $\mathcal{B}$'s advantage on that pair. Averaging over the draw of
the pair gives $\AdvH_{\Pi,\mathcal{A}} \ge \delta/2$, hence
$\delta \le 2\,\AdvH_{\Pi,\mathcal{A}}(\lambda)$. The averaging is necessary because $\mathcal{B}$'s
advantage may be concentrated on a small set of histories; the bound holds for the average, not
pointwise. This is the standard equivalence between semantic security and indistinguishability,
instantiated for the cohort-H game.

The direction of the corollary is worth restating, because it is easy to invert by accident. It
converts a \emph{proof} of dual-epoch security into a bound on predicate recovery. It does not
convert a \emph{measurement} into such a bound: an experiment exhibits one attacker, and the
advantage it achieves is a lower bound on the leakage present, never an upper bound on the leakage
available.

\paragraph{Necessity of the conditions} Condition S2 requires pseudorandom rather than merely
CPA-secure ciphertexts because free slots are uniform strings that must be indistinguishable from
live ones. Condition S3's whole-region rewrite is used in item 3: a per-slot write would reveal
which slot changed and, accumulated, the occupancy that item 2 hides. S3's padded refusal is used in
item 6 and is what lets the theorem cover overflow without an admissibility assumption; dropping it
for silent refusal reopens the capacity channel of \cref{prop:capacity}, which
\cref{sec:eval-refusal} measures. The flush at $\tau$ required by the game closes the
staged-remainder channel. \Cref{prop:seal,prop:fp} show that dropping S1, S2 or S3 admits a
distinguisher, which \cref{sec:eval-game} realizes as ablations, together with the observation that
the keyed relabeling is \emph{not} among the conditions and that removing it leaves the game at
chance.

\section{Proof of Theorem~\ref{thm:reactive}}
\label{supp:reactive}

We prove that a store satisfying S1, S2 and the clock conditions C1--C3 is reactively
$\mathcal{L}_{\mathrm{clock}}$-dual-epoch secure (\cref{def:reactive}) against every polynomial-time
feedback agent $\mathsf{P}$. The structure mirrors the proof of \cref{thm:main}; what changes is that
the operation sequence after $\tau$ is produced by $\mathsf{P}$ reacting to plaintext search results,
so the argument must show that this reaction never reaches the adversary's view.

\paragraph{Setup} Fix $\mathsf{P}$ and an adversary $\mathcal{A}$. In
$\mathbf{G}^{\mathrm{R}}_{\Pi,\mathsf{P},\mathcal{A}}$ the challenger builds cohort~H from $\sigma_b$,
flushes at $\tau$, and then for each adversary-supplied input runs one turn of $\mathsf{P}$: a
sequence of $\mathsf{search}$ operations whose plaintext results $\mathsf{P}$ reads, followed by
$\mathsf{put}$, $\mathsf{delete}$ or $\mathsf{update}$ operations that may depend on those results.
The adversary receives $\mathcal{L}^{\mathrm{op}}$ after each turn and, separately, the clock leakage
$\mathcal{L}_{\mathrm{clock}}$ of unlock, lock and flush times.

\paragraph{Game $\mathrm{H}_0$} The real reactive game.

\paragraph{Game $\mathrm{H}_1$} Replace every pseudorandom function under a store key by a lazily
sampled random function, as in $\mathrm{G}_1$. The gap is $\mathsf{Adv}^{\mathrm{prf}}_F(\lambda)$ up
to the number of derived keys. This step is oblivious to who chose the operations, so $\mathsf{P}$'s
reactivity does not affect it.

\paragraph{Game $\mathrm{H}_2$} Replace every ciphertext written to disk by a uniformly random string
of the same length (C2), keeping plaintexts in challenger state so no stored ciphertext is decrypted.
The gap is $\mathsf{Adv}^{\mathrm{ind\$}}_{\mathsf{AE}}(\lambda)$. After this step every byte on disk,
in every slot and in the sealed journal that carries staged records between flushes, is a uniform
string whose length is a public parameter.

\paragraph{The view is a function of the clock} We show the adversary's view in $\mathrm{H}_2$ is
determined by $\mathcal{L}_{\mathrm{clock}}$ and public randomness, independently of $b$, of the
inputs, and of every branch $\mathsf{P}$ took.

\begin{enumerate}
\item \emph{Snapshot.} By S1 and $\mathrm{H}_2$ the snapshot is a tuple of fixed-length uniform
strings whose count and lengths are public, identical under $b=0$ and $b=1$.
\item \emph{Writes.} By C1 the set of disk locations a flush rewrites is the image of the flush index
under the deterministic sweep: the $\kappa$ slots after a cursor that advances by $\kappa$ each flush from a
start fixed at unlock. This set is a function of the flush index alone. It does not depend on how many
records $\mathsf{P}$ staged, on their content, or on whether $\mathsf{P}$ wrote at all --- a flush
occurs every $\Delta$ regardless (C1). By $\mathrm{H}_2$ the bytes written are uniform. So the write
transcript is the clock's sequence of cursor ranges, filled with uniform bytes.
\item \emph{Reads.} By C3 a $\mathsf{search}$ issues no disk access after the unlock scan. Therefore
the $\mathsf{search}$ operations $\mathsf{P}$ performs, and the results it reads, produce \emph{no}
storage event. This is the crux: the only place $\mathsf{P}$'s secret-dependent behaviour could enter
the transcript is through the operations it issues, and its reads issue nothing while its writes are
placed by the clock, not by content. The branch $\mathsf{P}$ takes is thus invisible.
\item \emph{The journal.} Staged records wait in the fixed-size sealed journal, which C1 rewrites in
full every flush whether it holds staged records or none; by $\mathrm{H}_2$ its bytes are uniform of
fixed length. Its rewrites are on the clock and carry nothing.
\item \emph{Deletions and updates.} A deletion writes nothing at call time (C1 admits only the
clocked sweep as a writer); its slot is freed when the sweep reaches it, an event already accounted
in item~2. An update is a deletion and a staged write, both covered.
\end{enumerate}

Every component is a function of the unlock/lock/flush clock and uniform bytes. A simulator given
$\mathcal{L}_{\mathrm{clock}}$ emits the fixed-length uniform snapshot, and at each flush time emits
the cursor range that flush index determines, filled with uniform bytes; it runs neither $\mathsf{P}$
nor the challenge, and never sees $b$ or the inputs. Hence the real and simulated views are
indistinguishable and
\[
\AdvR_{\Pi,\mathsf{P},\mathcal{A}}(\lambda) \le \mathsf{Adv}^{\mathrm{prf}}_F(\lambda) +
\mathsf{Adv}^{\mathrm{ind\$}}_{\mathsf{AE}}(\lambda) \le \negl(\lambda).
\]
\hfill$\blacksquare$

\paragraph{The lock} C1 includes one flush at lock, always. A lock is a time the adversary chose
(\cref{def:reactive}), so its flush is a clock event like any other, and the cursor range it writes
is the next one in the sweep. A store that flushed \emph{more} than once at lock when more records
were staged than its journal holds would make the number of lock-time flushes a function of how much
$\mathsf{P}$ wrote, which violates C1; this is why Mnemos-M's back-pressure, which refuses a
$\mathsf{put}$ to the caller rather than stage beyond the journal, is the configuration
\cref{cor:m-reactive} covers without exception.

\paragraph{Part (i)} Let $\Pi$ be operation-oblivious and let $\mathsf{P}$ issue at most $q$ turns.
Define hybrids $\mathrm{T}_0,\dots,\mathrm{T}_q$, where in $\mathrm{T}_j$ the operations of the
first $j$ turns are replaced by the empty sequence at the same times and the remaining turns run
$\mathsf{P}$ as in the real game. Consecutive hybrids differ in the operations of one turn, and by
obliviousness the adversary's view of that turn, given everything before it, is indistinguishable
under the two; the later turns are a function of that view and of $\mathsf{P}$'s state, which the
adversary never sees. So $|\Pr[\mathrm{T}_{j-1}\Rightarrow 1] - \Pr[\mathrm{T}_j\Rightarrow 1]| \le
\negl(\lambda)$, and in $\mathrm{T}_q$ the view is the snapshot followed by an idle schedule, both
independent of $b$ by S1, S2 and the obliviousness of idle turns across histories. Summing, $\AdvR
\le q\cdot\negl(\lambda)$.

\paragraph{Part (ii)} Fix a store state and two operation sequences $O$, $O'$, searches included,
whose leakage is at statistical distance $\varepsilon$, and let $T$ be the optimal test between the
two leakage distributions, with advantage $\varepsilon$. Let $x^\star$ be a fact absent from that
state, let $\sigma_1$ produce the state and $\sigma_0$ produce it with $x^\star$ added, both padded to
equal length. $\mathsf{P}$, on input $x$, searches for $x$ and issues $O$ if the top result is a
near-duplicate of $x$ and $O'$ otherwise. $\mathcal{A}$ supplies $x^\star$ and outputs $b'=0$ exactly
when $T$ says $O$. Under $b=1$ no near-duplicate exists and $\mathsf{P}$ issues $O'$; under $b=0$ it
issues $O$ except with the probability $\delta$ that the search misses the stored duplicate. Two
further differences remain between the worlds: the leakage of $\mathsf{P}$'s search, and the
difference between the two store states on which $O$ and $O'$ run. An adversary of
\cref{def:dualepoch} can issue the same search and the same operations on either state itself, so by
dual-epoch security both differences are indistinguishable. Hence $\Pr[b'=b] \ge \tfrac12 +
\tfrac{\varepsilon}{2} - \delta - \negl(\lambda)$, and $\AdvR \ge \varepsilon/2 - \delta -
\negl(\lambda)$. \Cref{prop:derived} is the case where $O$ is a write and $O'$ is nothing.

\paragraph{Why Mnemos-P is excluded} The proof uses C1 and C3 exactly where \cref{prop:derived}'s
attack strikes. Mnemos-P flushes on demand, violating C1: a flush occurs \emph{because} $\mathsf{P}$
staged a record, so the occurrence of a flush is the write-occurrence bit a deduplicating agent
leaks. And Mnemos-P places by content, so even under a clocked flush the region a reactive write
lands in would depend on the retrieved plaintext. The clock-only store closes both, which is why the
composition holds for it and not for the partitioned design --- the separation
$\mathrm{RDE}\Rightarrow\mathrm{DE}$, $\mathrm{DE}\not\Rightarrow\mathrm{RDE}$ of
\cref{cor:m-reactive}. \Cref{sec:eval-reactive} measures both sides.
